% Section 1.2, from lectures/L01/README.md: the goals, the evaluation questions, and the next
% lecture.

\needspace{10\baselineskip}
\section{Review}
\label{c1:sec:review}

\subsection*{What you should be able to do now}
After this chapter, you should:
\begin{itemize}
  \item Be comfortable enough with C's core syntax (types, structs, functions, pointers) to not be
    fighting the language once hardware programming starts in the next chapter.
  \item Be able to define and use a \code{struct}, \code{enum}, and \code{typedef}, and explain why
    \code{self} is passed as a pointer.
  \item Be able to write a \code{.c}/\code{.h} pair with a correct header guard and a sensible
    static/public ordering.
  \item Understand the difference between a single pointer (\code{int*}) and a double pointer
    (\code{int**}) as a function parameter, and when each is needed.
  \item Be able to use \code{malloc}/\code{realloc}/\code{free} correctly, including why a pointer
    should be nulled immediately after \code{free}.
\end{itemize}

\needspace{6\baselineskip}
\subsection*{Questions to test yourself}
You should be able to explain:
\begin{enumerate}
  \item Why is \code{self} passed as a pointer rather than by value in this codebase's struct
    functions?
  \item What's the difference between a single pointer (\code{int*}) and a double pointer
    (\code{int**}) as a function parameter, and when do you need the second?
  \item Why must you capture \code{realloc}'s return value in a temporary before overwriting your
    original pointer?
  \item Why should you set a pointer to \code{0}/\code{NULL} immediately after \code{free}?
\end{enumerate}

\needspace{6\baselineskip}
\subsection*{Looking ahead}
Chapter~\ref{c2:ch} is about hardware-near programming in C: registers, bitwise I/O, and
driving/reading pins on the ATmega328P.
